% 作者题证的编辑转录副本，不是独立下载的上游原件。
% 来源：https://github.com/OpenLogicProject/OpenLogic/blob/master/content/incompleteness/incompleteness-provability/rosser-thm.tex

\paragraph{Notation (editor's expansion of the source macros).}
$\mathrm Q$ denotes Robinson arithmetic. For a formula $A$, $\ulcorner A\urcorner$ is the standard numeral of its G\"odel number; $\operatorname{diag}$ is the primitive recursive diagonal-substitution function. The formula $D_{\operatorname{diag}}(x,y)$ represents this function in $\mathrm Q$.
\begin{lemma}[Fixed-point lemma; the author's support proof]
Let $B(x)$ be any formula with one free variable $x$. Then there is a sentence $A$ such that $\mathrm Q\vdash A\leftrightarrow B(\ulcorner A\urcorner)$.
\end{lemma}
\begin{proof}
Given $B(x)$, let $E(x)$ be the formula $\exists y(D_{\operatorname{diag}}(x,y)\land B(y))$ and let $A$ be its diagonalization, i.e., the formula $E(\ulcorner E(x)\urcorner)$.
Since $D_{\operatorname{diag}}$ represents $\operatorname{diag}$, and $\operatorname{diag}(\ulcorner E(x)\urcorner)=\ulcorner A\urcorner$, $\mathrm Q$ can derive
\begin{align}
&D_{\operatorname{diag}}(\ulcorner E(x)\urcorner,\ulcorner A\urcorner),\label{repdiag1}\\
&\forall y(D_{\operatorname{diag}}(\ulcorner E(x)\urcorner,y)\to y=\ulcorner A\urcorner).\label{repdiag2}
\end{align}
Now we show that $\mathrm Q\vdash A\leftrightarrow B(\ulcorner A\urcorner)$. We argue informally, using just logic and facts derivable in $\mathrm Q$.
First, suppose $A$, i.e., $E(\ulcorner E(x)\urcorner)$. Going back to the definition of $E(x)$, we see that $E(\ulcorner E(x)\urcorner)$ just is
\[\exists y(D_{\operatorname{diag}}(\ulcorner E(x)\urcorner,y)\land B(y)).\]
Consider such a $y$. Since $D_{\operatorname{diag}}(\ulcorner E(x)\urcorner,y)$, by \eqref{repdiag2}, $y=\ulcorner A\urcorner$. So, from $B(y)$ we have $B(\ulcorner A\urcorner)$.
Now suppose $B(\ulcorner A\urcorner)$. By \eqref{repdiag1}, we have
\[D_{\operatorname{diag}}(\ulcorner E(x)\urcorner,\ulcorner A\urcorner)\land B(\ulcorner A\urcorner).\]
It follows that
\[\exists y(D_{\operatorname{diag}}(\ulcorner E(x)\urcorner,y)\land B(y)).\]
But that's just $E(\ulcorner E(x)\urcorner)$, i.e., $A$.
\end{proof}

\begin{theorem}[Rosser's theorem; the counted problem]
Let $T$ be any consistent, axiomatizable theory extending $\mathrm Q$. Then $T$ is not complete.
\end{theorem}
\begin{proof}
Recall that $\operatorname{Prov}_T(y)$ is defined as $\exists x\operatorname{Prf}_T(x,y)$, where $\operatorname{Prf}_T(x,y)$ represents the decidable relation which holds iff $x$ is the G\"odel number of a derivation of the sentence with G\"odel number $y$. The relation that holds between $x$ and $y$ if $x$ is the G\"odel number of a \emph{refutation} of the sentence with G\"odel number $y$ is also decidable. Let $\operatorname{not}(x)$ be the primitive recursive function which does the following: if $x$ is the code of a formula $A$, $\operatorname{not}(x)$ is a code of $\neg A$. Then $\operatorname{Refut}_T(x,y)$ holds iff $\operatorname{Prf}_T(x,\operatorname{not}(y))$. Let the corresponding arithmetic formula represent it. Then, if $T\vdash\neg A$ and $\delta$ is a corresponding derivation, $\mathrm Q\vdash\operatorname{Refut}_T(\ulcorner\delta\urcorner,\ulcorner A\urcorner)$. We define $\operatorname{RProv}_T(y)$ as
\[\exists x\bigl(\operatorname{Prf}_T(x,y)\land\forall z(z<x\to\neg\operatorname{Refut}_T(z,y))\bigr).\]
Roughly, $\operatorname{RProv}_T(y)$ says ``there is a proof of $y$ in $T$, and there is no shorter refutation of $y$.'' Assuming $T$ is consistent, $\operatorname{RProv}_T(y)$ is true of the same numbers as $\operatorname{Prov}_T(y)$; but from the point of view of \emph{provability} in $T$ (and we now know that there is a difference between truth and provability!) the two have different properties. If $T$ is \emph{in}consistent, then the two do \emph{not} hold of the same numbers!
By the fixed-point lemma, there is a formula $R_T$ such that
\begin{equation}\mathrm Q\vdash R_T\leftrightarrow\neg\operatorname{RProv}_T(\ulcorner R_T\urcorner).\label{RT}\end{equation}
In contrast to the proof of the first incompleteness theorem, here we claim that if $T$ is consistent, $T$ doesn't derive $R_T$, and $T$ also doesn't derive $\neg R_T$. (In other words, we don't need the assumption of $\omega$-consistency.)
First, let's show that $T\nvdash R_T$. Suppose it did, so there is a derivation of $R_T$ from $T$; let $n$ be its G\"odel number. Then $\mathrm Q\vdash\operatorname{Prf}_T(\bar n,\ulcorner R_T\urcorner)$, since $\operatorname{Prf}_T$ represents the proof relation in $\mathrm Q$. Also, for each $k<n$, $k$ is not the G\"odel number of a derivation of $\neg R_T$, since $T$ is consistent. So for each $k<n$, $\mathrm Q\vdash\neg\operatorname{Refut}_T(\bar k,\ulcorner R_T\urcorner)$. By the bounded-numeral lemma \textsf{lem:less-nsucc},
\[\mathrm Q\vdash\forall z(z<\bar n\to\neg\operatorname{Refut}_T(z,\ulcorner R_T\urcorner)).\]
Thus,
\[\mathrm Q\vdash\exists x\bigl(\operatorname{Prf}_T(x,\ulcorner R_T\urcorner)\land\forall z(z<x\to\neg\operatorname{Refut}_T(z,\ulcorner R_T\urcorner))\bigr),\]
but that's just $\operatorname{RProv}_T(\ulcorner R_T\urcorner)$. By \eqref{RT}, $\mathrm Q\vdash\neg R_T$. Since $T$ extends $\mathrm Q$, also $T\vdash\neg R_T$. We've assumed that $T\vdash R_T$, so $T$ would be inconsistent, contrary to the assumption of the theorem.
Now, let's show that $T\nvdash\neg R_T$. Again, suppose it did, and suppose $n$ is the G\"odel number of a derivation of $\neg R_T$. Then $\operatorname{Refut}_T(n,\ulcorner R_T\urcorner)$ holds, and since the corresponding formula represents it in $\mathrm Q$, $\mathrm Q\vdash\operatorname{Refut}_T(\bar n,\ulcorner R_T\urcorner)$. We'll again show that $T$ would then be inconsistent because it would also derive $R_T$. Since
\[\mathrm Q\vdash R_T\leftrightarrow\neg\operatorname{RProv}_T(\ulcorner R_T\urcorner),\]
and since $T$ extends $\mathrm Q$, it suffices to show that
\[\mathrm Q\vdash\neg\operatorname{RProv}_T(\ulcorner R_T\urcorner).\]
The sentence $\neg\operatorname{RProv}_T(\ulcorner R_T\urcorner)$, i.e.,
\[\neg\exists x\bigl(\operatorname{Prf}_T(x,\ulcorner R_T\urcorner)\land\forall z(z<x\to\neg\operatorname{Refut}_T(z,\ulcorner R_T\urcorner))\bigr),\]
is logically equivalent to
\[\forall x\bigl(\operatorname{Prf}_T(x,\ulcorner R_T\urcorner)\to\exists z(z<x\land\operatorname{Refut}_T(z,\ulcorner R_T\urcorner))\bigr).\]
We argue informally using logic, making use of facts about what $\mathrm Q$ derives. Suppose $x$ is arbitrary and $\operatorname{Prf}_T(x,\ulcorner R_T\urcorner)$. We already know that $T\nvdash R_T$, and so for every $k$, $\mathrm Q\vdash\neg\operatorname{Prf}_T(\bar k,\ulcorner R_T\urcorner)$. Thus, for every $k$ it follows that $x\ne\bar k$. In particular, we have (a) that $x\ne\bar n$. We also have $\neg(x=\bar0\lor x=\bar1\lor\cdots\lor x=\overline{n-1})$ and so by \textsf{lem:less-nsucc}, (b) $\neg(x<\bar n)$. By \textsf{lem:trichotomy}, $\bar n<x$. Since $\mathrm Q\vdash\operatorname{Refut}_T(\bar n,\ulcorner R_T\urcorner)$, we have $\bar n<x\land\operatorname{Refut}_T(\bar n,\ulcorner R_T\urcorner)$, and from that
\[\exists z(z<x\land\operatorname{Refut}_T(z,\ulcorner R_T\urcorner)).\]
Since $x$ was arbitrary we get, as required, that
\[\forall x\bigl(\operatorname{Prf}_T(x,\ulcorner R_T\urcorner)\to\exists z(z<x\land\operatorname{Refut}_T(z,\ulcorner R_T\urcorner))\bigr).\]
\end{proof}
